\documentclass[a4paper,12pt]{article}
\usepackage[T1]{fontenc}
\usepackage[utf8]{inputenc}
\usepackage[catalan]{babel}
\pagestyle{empty}
\usepackage{amsmath}
\usepackage{amssymb}
\usepackage{multicol}
\usepackage[most]{tcolorbox}
\usepackage{tikz}
\usetikzlibrary{decorations.pathmorphing}
\usetikzlibrary{arrows,positioning,shapes,fit,calc}
\setlength{\topmargin}{-0.3in} \setlength{\headheight}{0in}
\setlength{\headsep}{0in} \setlength{\footskip}{0in}
\setlength{\evensidemargin}{-.7in}
\setlength{\oddsidemargin}{-.7in} \setlength{\textheight}{9.3in}
\setlength{\textwidth}{7.5in} \setlength{\parindent}{0in}

\begin{document}
\textbf{2a Avaluació \hfill{}Física\hfill{}1r Batxillerat}\\
\textbf{Pla inclinat  \hfill{}\hfill{}Data:$\qquad\qquad$}\\
\textbf{Nom i cognoms:}\rule{2.4in}{.01in}
\hfill{}\textbf{Qualificació:}\rule{1in}{.01in} \vspace{0.5cm}

\small
\hspace{.25in} \textit{Instruccions:\/{}} Feu els exercicis a
l'espai que se us proporciona. Feu servir la cara posterior si
necessiteu més espai, \textit{indiqueu-ho clarament en aquest
cas.\/{}} Heu d'identificar clarament les respostes i mostrar el
procés per tal d'aconseguir la màxima puntuació. La puntuació dels
exercicis es dona entre parèntesis.
\\
\rule[.2in]{7.5in}{.01in}
\normalsize

\begin{enumerate}
\item
 Considereu el sistema dinàmic següent i siguin $m_{1}=m_{2}=20\,kg$ les masses que hi apareixen. Anomeneu $\alpha=60^{\circ}$ i considereu que el coeficient de fregament val $\mu=0,1$. Podeu suposar que el sistema es mou ``cap a l'esquerra'', degut a l'acció de la força $F$ que val $500\,N$. Es demana: 
\begin{enumerate}
\item 
\textbf{(2 pts)} Dibuixeu clarament les forces que actuen sobre cada massa.
\item
\textbf{(2 pts)} Escriviu les equacions (amb lletres) corresponents a cada massa que permeten resoldre el problema.
\item
\textbf{(1 pt)} Resoleu el sistema format per les equacions anteriors per trobar l'acceleració, en funció de $m_{1}$, $m_{2}$, $\alpha$, $\mu$, $F$ i $g$.
\item
\textbf{(1 pt)}
Finalment, trobeu el valor numèric de l'acceleració fent servir les dades proporcionades a l'enunciat. 
\end{enumerate}
\begin{center}
\begin{tikzpicture}[scale=1.5]
%\draw()--()--();
%\draw[step=0.5,gray!40!](-4,-4)grid(4,4);
%\filldraw(0,0)circle(0.05);
\begin{scope}[rotate around={29.5:(-3,0)}]
\draw(-2,0)rectangle(-1,1);
\filldraw(-1.5,0.5)circle(0.025);
%\draw[->](-1.5,0.5)--+(1,0);
%\draw[->](-1.5,0.5)--+(-1,0);
%\draw[->](-1.5,0.5)--+(-1.5,0);
%\draw[->](-1.5,0.5)--+(0,1);
%\draw[->](-1.5,0.5)--+(0,-1.8);
%\draw(-1.5,-0.5)arc(270:240:1);
%\draw[->,rotate around={240.5:(-1.5,0.5)}](-1.5,0.5)--+(2,0);
\draw(-3,0)--(1.04,0);
\draw[->](-2,0.5)--+(-1,0);
\draw(-1,0.5)--+(2.49,0);
\end{scope}
\node at (-3.5,0.75) {$F$};
\draw(0.75,2.5)circle(0.18);
\draw(0.5,2)--(0.5,-1);
\draw(0.93,2.5)--(0.93,1);
\draw(0.7,0)rectangle(1.16,1);
\draw(0.5,1)arc(270:211:1);
\node[scale=1] at (0,0.75) {$\alpha$};
%\node[scale=0.75] at (-1.7,-0.25) {$\beta$};
\node[scale=1] at (-1.75,2.25) {$m_{1}$};
%\node[scale=0.55] at (-1,-0.6) {$P_{1y}$};
%\node[scale=0.55] at (-2.9,0.95) {$P_{1x}$};
%\node[scale=0.55] at (-3.45,0.75) {$F_{f}$};
%\node[scale=0.75] at (-1,2) {$T$};
%\node[scale=0.75] at (1.25,1.6) {$T$};
%\node[scale=0.75] at (-2.5,2.25) {$N$};
\node[scale=1] at (1.5,-0.5) {$m_{2}$};
%\draw[->](0.93,0.5)--(0.93,-0.75);
%\draw[->](0.93,0.5)--(0.93,1.75);
\filldraw(0.93,0.5)circle(0.025);
\draw(0.5,2)--+(0.25,0.5);
\end{tikzpicture}
\end{center}
%\vspace{3.5cm}
\clearpage
\item
Considereu el sistema dinàmic següent i siguin $m_{1}=m_{2}=10\,kg$ les masses que hi apareixen. Anomeneu $\alpha=20^{\circ}$, $\beta=50^{\circ}$ i considereu que el coeficient de fregament val $\mu=0,1$. Podeu suposar que el sistema es mou ``cap a la dreta''. Es demana: 
\begin{enumerate}
\item 
\textbf{(2 pts)} Dibuixeu clarament les forces que actuen sobre cada massa.
\item
\textbf{(2 pts)} Escriviu les equacions (amb lletres) corresponents a cada massa que permeten resoldre el problema.
\item
\textbf{(1 pt)} Resoleu el sistema format per les equacions anteriors per trobar l'acceleració, en funció de $m_{1}$, $m_{2}$, $\alpha$, $\beta$, $\mu$ i $g$.
\item
\textbf{(1 pt)}
Finalment, trobeu el valor numèric de l'acceleració fent servir les dades proporcionades a l'enunciat. 
\end{enumerate}

\begin{center}
\begin{tikzpicture}
\draw(-5,0)--(7,0);
\begin{scope}[rotate around={45:(-5,0)}]
\draw(-5,0)--(1.2,0);
\draw(-3,0) rectangle (-1,2);
%\draw[->](-2,1)--(-2,3);
\begin{scope}[rotate around={-45:(-2,1)}]
%\draw[->](-2,1)--(-2,-2);
\end{scope}
%\draw[->,densely dashed](-2,1)--(-4.12,1);
%\draw[->](-1,1)--(0,1);
\draw(-2,1)--(1,1);
\draw(1,1)--(2,0.95);
%\draw[densely dashed,->](-2,1)--(-2,-1.12);
%\draw[densely dashed](-4,1)--(-5,1);
%\draw[loosely dashed] (-2,-1.12)--(-4.12,-1.12);
%\draw[loosely dashed](-4.12,-1.12)--(-4.12,1);
\end{scope}
\draw[gray](-0.6,4.4)--(-0.6,5.15);
\draw(-0.6,5.4) circle (0.25);
 \begin{scope}[rotate around={-30:(7,0)}]
 \draw(-1.8,0)--(7,0);
 \draw(1.5,0) rectangle (3.5,2);
 %\draw[->](2.5,1)--(2.5,3);
 \begin{scope}[rotate around={30:(2.5,1)}]
 %\draw[->](2.5,1)--(2.5,-2);
 \end{scope}
 %\draw[densely dashed](4,1)--(5,1);
 \draw(2.5,1)--(0,1);
 %\draw[->](4,1)--(5,1);
 \draw(0,1)--(-1,1);
 \draw(-1,1)--(-2.1,1.05);
%\draw[->,densely dashed](-2,1)--(1,1);
 %\draw[densely dashed,->](2.5,1)--(4,1);
 %\draw[densely dashed,->](2.5,1)--(2.5,-1.598);
%\draw[loosely dashed](4,1)--(4,-1.598);
%\draw[loosely dashed](4,-1.598)--(2.5,-1.598);
\end{scope}
%\node at (-2.5,4.5) {$F_{f1}$};
%\node at (-5.25,1.95) {$P_{1x}$};
%\node at (-5,4.5) {$N_{1}$};
\node at (-3.5,4.5) {$m_{1}$};
%\node at (-2,0.75) {$P_{1y}$};
%\node at (-1.5,5.45) {$T$};
%\node at (0.5,4.5) {$T$};
%\node at (5.75,1.5) {$F_{f2}$};
%\node at (1.75,1) {$P_{2y}$};
\node at (3.5,5) {$m_{2}$};
%\node at (5.25,2.8) {$P_{2x}$};
%\node at (5,5) {$N_{2}$};
\node at (-4.4,0.25) {$\alpha$};
\draw(-3.5,0)arc(0:45:1.5);
%\node at (-3.4,2.3) {$\alpha$};
\node at (5.5,0.26) {$\beta$};
\draw(5.75,0)arc(180:150:1.25);
%\node at (3.4,2.27) {$\beta$};
%\draw [<-] (-1.25,6) arc [start angle=160, end angle=20, radius=20pt];
%\node at (-0.5,7) {$\vec{a}$};
%\draw[step=1cm,gray,ultra thin] (-6,-1) grid (6,5);
%\filldraw (0,0) circle (0.1cm);
\end{tikzpicture}    
\end{center}
\clearpage
\item
Considereu el sistema dinàmic següent i siguin $m_{1}=20\,kg$, $m_{2}=15\,kg$ i $m_{3}=45\,kg$ les masses que hi apareixen. Anomeneu $\alpha=30^{\circ}$ i considereu que el coeficient de fregament val $\mu=0,1$. Podeu suposar que el sistema es mou ``cap a l'esquerra'' degut a l'acció de la força $F$, que val $800\,N$. Es demana: 
\begin{enumerate}
\item 
\textbf{(3 pts)} Dibuixeu clarament les forces que actuen sobre cada massa.
\item
\textbf{(3 pts)} Escriviu les equacions (amb lletres) corresponents a cada massa que permeten resoldre el problema.
\item
\textbf{(1 pt)} Resoleu el sistema format per les equacions anteriors per trobar l'acceleració, en funció de $F$, $m_{1}$, $m_{2}$, $m_{3}$, $\alpha$, $\mu$ i $g$.
\item
\textbf{(1 pt)}
Finalment, trobeu el valor numèric de l'acceleració fent servir les dades proporcionades a l'enunciat. 
\end{enumerate}
\begin{center}
\begin{tikzpicture}[scale=1.5]
%\draw()--()--();
%\draw[step=0.5,gray!40!](-7,-4)grid(4,4);
%\filldraw(0,0)circle(0.05);
\begin{scope}[rotate around={29.5:(-3,0)}]
\draw(-2,0)rectangle(-1,1);
\filldraw(-1.5,0.5)circle(0.025);
%\draw[->,red](-1.5,0.5)--+(1,0);
%\draw[->,green](-1.5,0.5)--+(-1,0);
%---------------------------------
%\draw[->,blue](-2.5,0.5)--(-3.4,0.5);
%\draw[red](-3.4,0.5)--(-3.86,0.5);
%----------------------------------
\draw(-2,0.5)--+(-1.85,0);
%\draw[->,orange](-1.5,0.5)--+(0,1);
%\draw[->,gray](-1.5,0.5)--+(0,-1.8);
%\draw(-1.5,-0.5)arc(270:240:1);
%\draw[->,rotate around={240.5:(-1.5,0.5)}](-1.5,0.5)--+(2,0);
\draw(-4,0)--(1.04,0);
%\draw(-1.5,-0.5)arc(270:240:1);
\draw(-1,0.5)--+(2.49,0);
\end{scope}
\draw[thick](0.75,2.5)circle(0.18);
\draw(0.5,2)--(0.5,-1);
\draw(0.93,2.5)--(0.93,1);
\draw(0.7,0)rectangle(1.16,1);

%\node[scale=0.75,green] at (-1.75,0.5) {$\alpha$};
\node[scale=0.75] at (-1.5,2.25) {$m_{2}$};
%\node[scale=0.55] at (-1,-0.6) {$P_{2y}$};
%\node[scale=0.55] at (-2.9,0.95) {$P_{2x}$};
%\node[scale=0.55] at (-3.45,0.75) {$F_{f2}$};
%\node[scale=0.75] at (-1,2) {$T'$};
%\node[scale=0.75] at (1.25,1.6) {$T'$};
%\node[scale=0.75] at (-2.5,2.25) {$N_{2}$};
\node[scale=0.75] at (1.5,-0.25) {$m_{3}$};
%\draw[->,green](0.93,0.5)--(0.93,-0.75);
%\draw[->,blue](0.93,0.5)--(0.93,1.75);
\filldraw(0.93,0.5)circle(0.025);
\draw(0.5,2)--+(0.25,0.5);
%\node[scale=0.75] at (-3.5,-0.1) {$T$};
%\node[scale=0.75] at (-4.5,-0.2) {$T$};
\draw(-4.1,0.25)--(-3.85,-0.5);
%----------------------
\draw(-2.5,-0.5)--(-7,-0.5);
\draw[thick](-4.1,0.25)circle(0.25);
\draw(-4.1,0)--(-5,0);
\draw(-5,0)--(-4.5,0);
\draw(-6,-0.5)rectangle(-5,0.5);
\draw[->](-6,0)--(-7,0);
\node at (-6.75,0.5) {$F$};
\filldraw(-5.5,0)circle(0.025);
%\draw[->,blue](-5.5,0)--+(-1,0);
%\draw[->,green](-5.5,0)--+(0,-1);
%\draw[->,red](-5.5,0)--+(0,1);
\draw(-5.5,0)--+(0.5,0);
\node[scale=0.75] at (-6,1) {$m_{1}$};
%\node[scale=0.75] at (-7,-0.25) {$F_{f1}$};
%\node[scale=0.75] at (-5.5,1.25) {$N_{1}$};
\draw(-3,-0.5)arc(0:34:0.75);
\node[scale=0.75] at (-2.75,-0.25) {$\alpha$};
\end{tikzpicture}
\end{center}
\end{enumerate}
\end{document}
